\documentclass[a4paper,11pt]{article}

% ---------- Sprache & Kodierung ----------
% Kompilierung mit XeLaTeX oder LuaLaTeX (Unicode-nativ, keine
% inputenc/fontenc-Pakete noetig).
\usepackage{fontspec}
\usepackage{polyglossia}
\setmainlanguage{german}

% ---------- Layout ----------
\usepackage[margin=2.5cm]{geometry}
\usepackage{parskip}
\usepackage{fancyhdr}
\usepackage{titling}
\usepackage{hyperref}

% ---------- Code-Listings ----------
\usepackage{xcolor}
\usepackage{listings}

\definecolor{codebg}{rgb}{0.96,0.96,0.96}
\definecolor{codegreen}{rgb}{0,0.5,0}
\definecolor{codegray}{rgb}{0.5,0.5,0.5}
\definecolor{codeblue}{rgb}{0,0,0.6}
\definecolor{framegray}{rgb}{0.8,0.8,0.8}

\lstdefinestyle{cppstyle}{
    backgroundcolor=\color{codebg},
    commentstyle=\color{codegreen},
    keywordstyle=\color{codeblue}\bfseries,
    numberstyle=\tiny\color{codegray},
    stringstyle=\color{codegreen},
    basicstyle=\ttfamily\footnotesize,
    breakatwhitespace=false,
    breaklines=true,
    captionpos=b,
    keepspaces=true,
    numbers=left,
    numbersep=6pt,
    showspaces=false,
    showstringspaces=false,
    showtabs=false,
    tabsize=4,
    frame=single,
    rulecolor=\color{framegray},
    language=C++,
}
\lstset{style=cppstyle}

\lstdefinestyle{cmakestyle}{
    backgroundcolor=\color{codebg},
    commentstyle=\color{codegreen},
    keywordstyle=\color{codeblue}\bfseries,
    numberstyle=\tiny\color{codegray},
    basicstyle=\ttfamily\footnotesize,
    breaklines=true,
    numbers=left,
    numbersep=6pt,
    frame=single,
    rulecolor=\color{framegray},
    tabsize=2,
}

% ---------- Kopf-/Fusszeile ----------
\setlength{\headheight}{14.5pt}
\addtolength{\topmargin}{-2.5pt}
\pagestyle{fancy}
\fancyhf{}
\fancyhead[L]{Chat-Server -- Codedokumentation}
\fancyhead[R]{\thepage}
\renewcommand{\headrulewidth}{0.4pt}

% ---------- Titeldaten ----------
\title{TCP-Chat-Server in C++\\\large Codedokumentation für die Fachprüfung}
\author{Christine Münzberg}
\date{28. Juli 2026}

\begin{document}

\maketitle
\thispagestyle{empty}

\tableofcontents
\newpage

% ============================================================
\section{Projektübersicht}

Dieses Projekt implementiert einen einfachen 1:1-Chat über TCP-Sockets unter
Windows (Winsock2-API). Das Projekt besteht aus zwei eigenständigen
Programmen:

\begin{itemize}
    \item \textbf{server.cpp} -- wartet auf eine eingehende Verbindung und
    tauscht anschließend abwechselnd Nachrichten mit dem Client aus.
    \item \textbf{client.cpp} -- baut eine Verbindung zum Server auf
    (Standard: \texttt{127.0.0.1:9000}) und kommuniziert im gleichen
    Frage-Antwort-Rhythmus.
\end{itemize}

Beide Programme werden unabhängig voneinander kompiliert und als separate
ausführbare Dateien (\texttt{server.exe}, \texttt{client.exe}) gestartet.
Der Build erfolgt entweder direkt mit \texttt{g++} oder über
\texttt{CMakeLists.txt}.

% ============================================================
\section{Architektur und Ablauf}

Beide Seiten folgen dem klassischen Winsock-Lebenszyklus. Der Unterschied
liegt darin, dass der Server passiv auf eine Verbindung wartet (\texttt{bind}
$\rightarrow$ \texttt{listen} $\rightarrow$ \texttt{accept}), während der
Client aktiv eine Verbindung aufbaut (\texttt{connect}).

\begin{center}
\begin{tabular}{|l|l|}
\hline
\textbf{Server} & \textbf{Client} \\
\hline
WSAStartup & WSAStartup \\
socket & socket \\
bind & -- \\
listen & -- \\
accept & connect \\
recv / send (Schleife) & send / recv (Schleife) \\
closesocket / WSACleanup & closesocket / WSACleanup \\
\hline
\end{tabular}
\end{center}

Nach dem Verbindungsaufbau laufen Server und Client in einer Schleife, in
der abwechselnd eine Nachricht empfangen und eine Nachricht gesendet wird.
Die Eingabe \texttt{"exit"} auf einer der beiden Seiten beendet den Chat
kontrolliert auf beiden Seiten.

\subsection{Wichtige Konzepte}

\begin{description}
    \item[TCP vs. UDP] TCP ist verbindungsorientiert und garantiert
    zuverlässige, geordnete Zustellung -- passend für einen Chat, bei dem
    keine Nachricht verloren gehen soll. UDP wäre verbindungslos und
    ungeordnet, aber schneller.
    \item[Blockierende Aufrufe] \texttt{accept()} und \texttt{recv()}
    blockieren das Programm, bis ein Ereignis eintritt (neue Verbindung
    bzw. neue Daten). Dadurch läuft der Code synchron ohne aktives Polling.
    \item[Zweck des Puffers] \texttt{char puffer[PUFFER]} nimmt die
    empfangenen Rohdaten (Bytes) auf. Die Größe (1024 Byte) begrenzt die
    maximale Nachrichtenlänge pro \texttt{recv()}-Aufruf.
    \item[Umwandlung char* $\rightarrow$ std::string] Nach dem Empfang wird
    der Puffer an der empfangenen Byte-Anzahl mit \texttt{'\textbackslash 0'}
    terminiert und anschließend in ein \texttt{std::string} überführt, um
    komfortabel mit C++-Stringfunktionen weiterzuarbeiten.
\end{description}

% ============================================================
\section{Bauanleitung}

\subsection{Variante A: direkt mit g++}
\begin{lstlisting}[style=cmakestyle,language=bash,numbers=none]
g++ -o server.exe server.cpp -lws2_32 -std=c++17
g++ -o client.exe client.cpp -lws2_32 -std=c++17
\end{lstlisting}

\subsection{Variante B: über CMake}
\begin{lstlisting}[style=cmakestyle,language=bash,numbers=none]
mkdir build
cd build
cmake ..
cmake --build .
\end{lstlisting}

Der Linker-Flag \texttt{-lws2\_32} bindet die Windows-Socket-Bibliothek ein
und ist zwingend nötig, da Winsock2 (im Gegensatz zu den POSIX-Sockets unter
Linux) eine separate Bibliothek ist, die nicht automatisch verlinkt wird.

% ============================================================
\section{Quellcode: server.cpp}

Der Server initialisiert Winsock, erstellt einen Socket, bindet ihn an
Port~9000, wartet auf eine Verbindung und tauscht danach Nachrichten mit
dem Client aus, bis eine Seite \texttt{"exit"} sendet.

\lstinputlisting[style=cppstyle,caption={server.cpp -- vollständiger Quellcode}]{server.cpp}

% ============================================================
\section{Quellcode: client.cpp}

Der Client baut eine Verbindung zum Server auf (Loopback-Adresse
\texttt{127.0.0.1}, Port~9000) und sendet als Erster eine Nachricht, bevor
er auf die Antwort des Servers wartet.

\lstinputlisting[style=cppstyle,caption={client.cpp -- vollständiger Quellcode}]{client.cpp}

% ============================================================
\section{Build-Konfiguration: CMakeLists.txt}

Definiert zwei separate Targets (\texttt{server}, \texttt{client}) und
verlinkt unter Windows jeweils die Winsock2-Bibliothek.

\lstinputlisting[style=cmakestyle,language=bash,caption={CMakeLists.txt}]{CMakeLists.txt}

% ============================================================
\section{Fehlerbehandlung im Code}

Beide Programme prüfen nach jedem kritischen Winsock-Aufruf
(\texttt{WSAStartup}, \texttt{socket}, \texttt{bind}, \texttt{listen},
\texttt{accept}, \texttt{connect}, \texttt{send}) den Rückgabewert und
brechen bei Fehlern mit einer Meldung auf \texttt{std::cerr} sowie
kontrolliertem Aufräumen (\texttt{closesocket}, \texttt{WSACleanup}) ab.

\begin{center}
\begin{tabular}{|p{4.5cm}|p{4.5cm}|p{4.5cm}|}
\hline
\textbf{Fehlermeldung} & \textbf{Ursache} & \textbf{Lösung} \\
\hline
\texttt{'g++' is not recognized} & PATH nicht korrekt gesetzt & Neues
Terminal öffnen, PATH prüfen \\
\hline
\texttt{undefined reference to WSAStartup} & \texttt{-lws2\_32} fehlt beim
Kompilieren & Linker-Flag ergänzen \\
\hline
\texttt{winsock2.h: No such file or directory} & Falscher/kein
Windows-Compiler & MinGW-w64-Toolchain installieren \\
\hline
Bind fehlgeschlagen & Port 9000 bereits belegt & Alten Serverprozess
beenden, erneut starten \\
\hline
\end{tabular}
\end{center}

% ============================================================
\section{Mögliche Erweiterungen}

\begin{itemize}
    \item Mehrere Clients gleichzeitig (ein Thread pro Verbindung)
    \item Zeitstempel und Absendername je Nachricht
    \item Broadcast an alle verbundenen Clients statt 1:1
    \item Robustere Fehlerbehandlung, z.\,B. Timeout bei \texttt{recv}
\end{itemize}

\end{document}
